\documentclass[12pt,letterpaper]{article}

% ===================== PAQUETES =====================
\usepackage[spanish,es-tabla]{babel}
\usepackage[utf8]{inputenc}
\usepackage[T1]{fontenc}
\usepackage{geometry}
\usepackage{graphicx}
\usepackage{float}
\usepackage{booktabs}
\usepackage{array}
\usepackage{longtable}
\usepackage{tabularx}
\usepackage{enumitem}
\usepackage{hyperref}
\usepackage{xcolor}
\usepackage{titlesec}
\usepackage{fancyhdr}
\usepackage{pgfgantt}
\usepackage{tikz}
\usepackage{pgfplots}

\pgfplotsset{compat=1.18}
\geometry{margin=2.5cm}
\hypersetup{colorlinks=true, linkcolor=black, urlcolor=blue, citecolor=black}

% ===================== DATOS =====================
\newcommand{\nombreProyecto}{VitalCode: Entorno Virtual de Primeros Auxilios Psicológicos}
\newcommand{\nombreCurso}{Ingeniería y Desarrollo Sustentable/Sostenible}
\newcommand{\docentes}{Pía Del Canto C. -- Aracelly Aravena F.}
\newcommand{\ayudantes}{Anyel Cortés -- Jorge Herrera}
\newcommand{\grupo}{VitalCode}
\newcommand{\integrantes}{Elliot Bravo, Catalina Rojas, Matías Collao, Benjamín Cortés}
\newcommand{\fecha}{Mayo 2026}

\pagestyle{fancy}
\fancyhf{}
\lhead{\grupo}
\rhead{Informe de Avance IDS}
\cfoot{\thepage}

\titleformat{\section}{\large\bfseries}{\thesection.}{0.5em}{}
\titleformat{\subsection}{\normalsize\bfseries}{\thesubsection.}{0.5em}{}

\begin{document}

% ===================== PORTADA =====================
\begin{titlepage}
    \centering
    \vspace*{1cm}
    {\Large \textbf{\nombreCurso}}\\[0.4cm]
    {\large Primer Entregable -- Informe de Avance}\\[1.5cm]
    {\Huge \textbf{\nombreProyecto}}\\[1cm]
    {\Large \textbf{Grupo:} \grupo}\\[0.4cm]
    {\large \textbf{Integrantes:} \integrantes}\\[0.4cm]
    {\large \textbf{Docentes:} \docentes}\\[0.2cm]
    {\large \textbf{Ayudantes:} \ayudantes}\\[1cm]
    {\large \textbf{ODS principal:} ODS 3 -- Salud y Bienestar}\\[2cm]
    {\large \fecha}
    \vfill
\end{titlepage}

\tableofcontents
\newpage

% ===================== RESUMEN =====================
\section*{Resumen}
\addcontentsline{toc}{section}{Resumen}
El presente informe desarrolla la propuesta \textit{VitalCode}, una aplicación orientada a la salud mental adolescente mediante recursos digitales de primeros auxilios psicológicos. El proyecto se vincula con el ODS 3 (Salud y Bienestar) para reducir barreras de acceso y estigma. Es una solución digital de bajo impacto material, planteada como alternativa sostenible frente a infraestructuras presenciales.

\textbf{Palabras clave:} salud mental, ODS 3, aplicación, entorno virtual, primeros auxilios psicológicos, sostenibilidad.

\newpage

% ===================== 1 CONTEXTO =====================
\section{Contexto del proyecto}

\subsection{Descripción general de la empresa}
\textbf{VitalCode} es una empresa de base tecnológica social para el bienestar psicológico adolescente. El producto es una aplicación con entorno virtual (Unity o Roblox Studio) donde usuarios acceden anónimamente a recursos de apoyo emocional. Es una solución preventiva y educativa que no reemplaza la atención profesional.

\subsection{Estructura organizacional}
\begin{table}[H]
\centering
\caption{Estructura organizacional propuesta para VitalCode}
\begin{tabularx}{\textwidth}{p{4cm} X p{2.2cm}}
\toprule
\textbf{Cargo / área} & \textbf{Función principal} & \textbf{Cantidad} \\
\midrule
Dirección general & Coordinar estrategia y alianzas. & 1 \\
Jefatura de proyecto & Planificación y presupuesto. & 1 \\
Psicólogo/a asesor/a & Validar contenidos y protocolos. & 1 \\
Especialista salud mental & Pertinencia ética. & 1 \\
Desarrolladores/as & Programación de la App. & 2 \\
Diseñador/a UX/UI & Accesibilidad y seguridad. & 1 \\
Ciberseguridad & Privacidad de datos. & 1 \\
Estudiantes en práctica & Apoyo y documentación. & 2 \\
\bottomrule
\end{tabularx}
\end{table}

% ===================== 2 PROBLEMA =====================
\section{Descripción del problema u oportunidad}
La salud mental adolescente presenta barreras como el temor al juicio y falta de recursos tempranos. El ODS 3 busca garantizar vida sana; considerando que más de 700.000 personas mueren por suicidio anualmente, se prioriza la prevención digital.

% ===================== 3 ANALISIS ENTORNO =====================
\section{Análisis del entorno}

\subsection{Análisis FODA}
\begin{table}[H]
\centering
\caption{Análisis FODA del proyecto VitalCode}
\begin{tabularx}{\textwidth}{p{3cm} X}
\toprule
\textbf{Dimensión} & \textbf{Elementos principales} \\
\midrule
Fortalezas & Bajo costo, anonimato, alineación ODS 3. \\
Oportunidades & Escalabilidad, alianzas educativas. \\
Debilidades & Dependencia de conectividad, no es terapia clínica. \\
Amenazas & Regulaciones de datos, brecha digital. \\
\bottomrule
\end{tabularx}
\end{table}

% ===================== 4 PLANEACION =====================
\section{Planeación administrativa del proyecto}

\subsection{Presupuesto estimado}
\begin{table}[H]
\centering
\caption{Presupuesto preliminar del proyecto}
\begin{tabularx}{\textwidth}{X r}
\toprule
\textbf{Ítem} & \textbf{Costo estimado (CLP)} \\
\midrule
Licencias y Gráficos & 30.000 \\
Dominio y Hosting & 25.000 \\
Difusión Encuesta & 10.000 \\
Asesoría Salud Mental & 80.000 \\
Imprevistos & 30.000 \\
\midrule
\textbf{Total estimado} & \textbf{175.000} \\
\bottomrule
\end{tabularx}
\end{table}

% --- GRÁFICO 1: DISTRIBUCIÓN DEL PRESUPUESTO ---
\begin{figure}[H]
\centering
\begin{tikzpicture}
\begin{axis}[
    ybar,
    bar width=20pt,
    ylabel={Costo en CLP},
    symbolic x coords={Licencias, Hosting, Difusión, Asesoría, Otros},
    xtick=data,
    nodes near coords,
    ymin=0, ymax=100000,
    width=0.8\textwidth,
    height=6cm,
    axis x line*=bottom,
    axis y line*=left,
    enlarge x limits=0.15,
]
\addplot[fill=blue!40] coordinates {(Licencias,30000) (Hosting,25000) (Difusión,10000) (Asesoría,80000) (Otros,30000)};
\end{axis}
\end{tikzpicture}
\caption{Distribución visual del presupuesto estimado.}
\end{figure}

\subsection{Cronograma preliminar}
\begin{figure}[H]
\centering
\begin{ganttchart}[
    hgrid, vgrid, x unit=0.8cm, y unit chart=0.55cm,
    title height=1, bar height=0.6
]{1}{8}
\gantttitle{Semanas}{8} \\
\gantttitlelist{1,...,8}{1} \\
\ganttbar{Encuesta}{1}{2} \\
\ganttbar{Diseño UX}{2}{3} \\
\ganttbar{Prototipo}{3}{6} \\
\ganttbar{Informe}{7}{8}
\end{ganttchart}
\caption{Carta Gantt preliminar.}
\end{figure}

% ===================== 5 ENCUESTA =====================
\section{Encuesta}

\subsection{Resultados preliminares}
\begin{table}[H]
\centering
\caption{Resultados cuantitativos (N=25)}
\begin{tabularx}{\textwidth}{X c}
\toprule
\textbf{Indicador} & \textbf{Resultado (\%)} \\
\midrule
Interés en App anónima & 88\% \\
Importancia del anonimato & 92\% \\
Uso de ejercicios de respiración & 75\% \\
Preferencia por recursos educativos & 64\% \\
\bottomrule
\end{tabularx}
\end{table}

% --- GRÁFICO 2: INTERÉS EN FUNCIONALIDADES ---
\begin{figure}[H]
\centering
\begin{tikzpicture}
\begin{axis}[
    xbar,
    xlabel={Porcentaje de Aceptación (\%)},
    symbolic y coords={Educación, Respiración, Anonimato, Interés Gral.},
    ytick=data,
    nodes near coords,
    nodes near coords align={horizontal},
    xmin=0, xmax=100,
    width=0.8\textwidth,
    height=6cm,
    enlarge y limits=0.2,
]
\addplot[fill=green!40] coordinates {(64,Educación) (75,Respiración) (92,Anonimato) (88,Interés Gral.)};
\end{axis}
\end{tikzpicture}
\caption{Nivel de interés por funcionalidad según encuesta.}
\end{figure}

% ===================== CONCLUSIONES =====================
\section{Conclusiones preliminares}
VitalCode responde a la necesidad de salud mental preventiva. El éxito del proyecto depende de la validación profesional de los contenidos y del análisis de la encuesta para ajustar el prototipo funcional en Unity/Roblox.

% ===================== REFERENCIAS =====================
\newpage
\section*{Referencias}
\addcontentsline{toc}{section}{Referencias}
Organización Mundial de la Salud. (2019). \textit{Suicide worldwide in 2019}. OMS. \\
Organización de las Naciones Unidas. (s. f.). \textit{Objetivo 3: Salud y Bienestar}. \\
Universidad Católica del Norte. (2026). \textit{Ingeniería y Desarrollo Sustentable}.

\end{document}
